\chapter{Exponential and Logarithm}\label{ch:exponential-logarithm}

Repeated growth multiplies its successive growth factors. Logarithm measures
relative change by accumulating the reciprocal. This chapter constructs those
computations and proves that they describe inverse functions. Its purpose is
also practical: later applications may choose the computation that suits them,
without rebuilding the theory for each choice.

Every input below is a valid represented real or complex number, not merely a
rational argument. Equality of values means \texttt{RealRaw.Equiv} or
\texttt{ComplexRaw.Equiv}, written $\simeq$. Algorithms operate on rational
intervals or boxes, and every refinement has justified convergence. Neither
Mathlib's real numbers nor its complex numbers enter these constructions.

\section{The exponential power series}\label{sec:exponential-power-series}

For a finite prefix, put
\[
 E_N(z)=\sum_{k=0}^{N}\frac{z^k}{k!}.
\]
Choose $R$ bounding $|\Re z|+|\Im z|$ (or $|x|$ for a real input).
If $N+2\geq 2R$, successive omitted
absolute terms have ratio at most $1/2$. The remaining finite prefixes
therefore lie within a geometric majorant of size
\[
 2\frac{R^{N+1}}{(N+1)!}.
\]
The argument concerns all later finite prefixes; it does not presume a limit
in an abstract completed field. For represented inputs we also refine the
input box until its contribution to the evaluation error is small enough.
Intersecting compatible enclosures produces a valid computation $E(z)$.

\lean{ComputableAnalysis.FiniteExponentialTaylor.scheduled_expTaylorPrefix_enclosure}
\lean{ComputableAnalysis.ModularForms.entireExponentialValue}
\lean{ComputableAnalysis.ExponentialComputations.representedTaylor_agrees}
\leanok

Finite coefficient identities, together with the tail bounds, prove
\[
 E(0)\simeq 1,\qquad E(z+w)\simeq E(z)E(w),\qquad
 E(z)E(-z)\simeq 1.
\]
Termwise differentiation is justified by a quantitative remainder, giving an
actual holomorphic derivative $E'=E$ on every local chart. The derivative is
an exact identity for arbitrary represented inputs; a formal coefficient
shift alone would not suffice.

\lean{ComputableAnalysis.ModularForms.entireExponential_addition}
\lean{ComputableAnalysis.ModularForms.entireExponential_derivative_value}
\leanok

Restricting this computation to the real axis defines $\exp(x)$. The imaginary
coordinate is proved zero. Positivity holds for every real input, including
negative inputs: the squared magnitude of $E(x/2)$ equals $E(x)$, and the
reciprocal identity excludes zero. The real restriction is injective.

\begin{theorem}[Real exponential]\label{thm:represented-real-exponential}
\lean{ComputableAnalysis.ExponentialComputations.exp_real_embedding}
\lean{ComputableAnalysis.ExponentialComputations.exp_positive}
\lean{ComputableAnalysis.ExponentialComputations.exp_injective}
\lean{ComputableAnalysis.ExponentialComputations.exp_congr}
\leanok
The computation $\exp$ is positive, representation invariant, and injective
on all valid represented real inputs.
\end{theorem}

\section{Compound interest computes the same function}\label{sec:compound-exponential}

The alternative suggested by repeated small growth is
\[
 C_n(z)=\left(1+\frac{z}{n+1}\right)^{n+1}.
\]
These are literal finite powers. Their implementation does not substitute a
power-series value. Negative and complex inputs are allowed; a positive base
is only needed for the usual financial interpretation.

The checked comparison uses the local logarithm power series. Its complex
error bounds control both coordinates: write
\[
 \|z\|_{\square}=\max(|\Re z|,|\Im z|).
\]
On its small chart, its constant coefficient is zero, its linear coefficient is one, and
its tail gives
\[
 \left\|\log(1+w)-w\right\|_{\square}\leq 16H^2\qquad
 (\|w\|_{\square}\leq H).
\]
For sufficiently large $n$, take $w=z/(n+1)$. Multiplying the error by
$n+1$ shows that $(n+1)\log(1+w)$ differs from $z$ by a computable bound of
order $1/(n+1)$. The local inverse identity, the exponential addition law,
and a bounded-chart continuity estimate then give
\[
 \left\|C_n(z)-E(z)\right\|_{\square}\leq \frac{K_z}{n+1}.
\]
Here $K_z$ and an initial index are computed from a rational bound on the
input box. No caller chooses an internal schedule.

\lean{ComputableAnalysis.ExponentialComputations.localLogarithm_linear_error}
\lean{ComputableAnalysis.ExponentialComputations.compoundApproximation_error}
\leanok

\begin{theorem}[Compound-interest limit]\label{thm:compound-function-agreement}
\lean{ComputableAnalysis.ExponentialComputations.compoundValue}
\lean{ComputableAnalysis.ExponentialComputations.compoundValue_equiv_powerSeries}
\lean{ComputableAnalysis.ExponentialComputations.compoundValue_congr}
\leanok
The executable stabilized limit of the finite compound-interest powers
agrees with $E(z)$ for every valid represented complex input, and hence with
$\exp(x)$ on the real axis.
\end{theorem}

The current error constants are deliberately generous. This proves a usable
mathematical interchange law, rather than a claim that compound interest is
the fastest numerical evaluator.

\section{A logarithm on the whole positive axis}\label{sec:positive-logarithm}

A local logarithm does not by itself cover all positive inputs. The existing
continuation construction produces a complex logarithm of any nonzero
represented complex input. For a positive real input $x$, take the real
coordinate of that constructed logarithm and call it $\log x$.

This does not impose a global single-valued complex logarithm. Different
continuation routes may have different imaginary coordinates, but the
checked exponential fiber theorem proves that their real coordinates agree.
Consequently the real logarithm is independent of the input representation
and of the auxiliary branch. Exponential positivity and its real-axis
injectivity give both inverse identities.

\begin{theorem}[Global inverse laws]\label{thm:global-exp-log-inverses}
\lean{ComputableAnalysis.ExponentialComputations.log_congr}
\lean{ComputableAnalysis.ExponentialComputations.log_exp}
\lean{ComputableAnalysis.ExponentialComputations.exp_log}
\leanok
For every valid represented real $a$ and every positive valid represented
real $x$,
\[
 \log(\exp a)\simeq a,\qquad \exp(\log x)\simeq x.
\]
\end{theorem}

The logarithm also has an actual local holomorphic extension at each positive
center. Its derivative there is the constructed reciprocal. In particular,
restricting this chart to positive real inputs gives the familiar law
\[
 \frac{d}{dx}\log x\simeq\frac1x.
\]
The local expansion has a proved quadratic remainder, which will identify an
independently computed reciprocal integral.

\lean{ComputableAnalysis.ExponentialComputations.logChart_real_agreement}
\lean{ComputableAnalysis.ExponentialComputations.logChart_derivative}
\lean{ComputableAnalysis.ExponentialComputations.log_increment_error}
\leanok

\section{Computing the reciprocal integral}\label{sec:integral-logarithm}

For an arbitrary positive represented endpoint $x$, use the fixed rational
parameter interval $[0,1]$ and the segment
\[
 s_x(t)=1+t(x-1),\qquad
 g_x(t)=\frac{x-1}{1+t(x-1)}.
\]
Define the normalized reciprocal integral by the actual rectangle-sum
computation
\[
 L(x)=\int_0^1 g_x(t)\,dt.
\]
This is the oriented meaning of $\int_1^x du/u$ used here. When $x<1$, the
factor $x-1$ supplies the reversed orientation automatically. The parameter
interval stays ordered and its entire image stays positive.

To construct bounds, search the valid input intervals until a positive lower
endpoint is found. Its minimum with $1$ gives a rational $\ell>0$ such that
$s_x(t)\geq\ell$ throughout the segment. A rational bound $A$ on $|x-1|$
then supplies a Lipschitz bound for $g_x$; the implementation uses the safe
constant
\[
 V=\frac{8A^2}{\ell^2}.
\]
Choose a finite uniform grid fine enough that its variation contribution is
within the requested error. At each grid point refine the integrand box to
control its evaluation width. Pad this box by the variation bound on that
cell. The resulting lower and upper rectangle sums have arbitrarily small
gap. Prefix intersections make them a valid represented value.

\lean{ComputableAnalysis.ExponentialComputations.segmentLower}
\lean{ComputableAnalysis.ExponentialComputations.positiveSegment}
\lean{ComputableAnalysis.ExponentialComputations.segmentIntegrand_variation}
\lean{ComputableAnalysis.ExponentialComputations.lipschitzBounds_gap}
\lean{ComputableAnalysis.ExponentialComputations.integralLog}
\leanok

The runtime reads the reciprocal samples and their bounds. It does not read
$\log x$ as its numerical answer. Validity of the resulting real would still
not certify its role as an integral; that is the next theorem.

\section{Why this integral is logarithm}\label{sec:integral-logarithm-proof}

Put $F_x(t)=\log(s_x(t))$. The local logarithm remainder and the segment's
positive lower bound give a uniform small-step estimate
\[
 |F_x(v)-F_x(u)-(v-u)g_x(u)|\leq K(v-u)^2.
\]
One may use $K=64A^2/\ell^2$ and a justified positive step bound. Every point
of every subdivision has the required domain evidence.

Subdivide a cell into $N$ equal pieces and telescope. The accumulated error
is at most $K(v-u)^2/N$. Letting this explicit rational error shrink shows
that $F_x(v)-F_x(u)$ lies between the cell's lower and upper derivative
rectangles. Applying the finite FTC and the tight reciprocal bounds proves
actual \texttt{Integral.HasIntegral} evidence. Since $\log 1\simeq0$, the
endpoint difference is $\log x$.

\begin{theorem}[Reciprocal integral identification]\label{thm:global-reciprocal-integral}
\lean{ComputableAnalysis.ExponentialComputations.cellBounds_of_quadratic}
\lean{ComputableAnalysis.ExponentialComputations.segmentLog_hasIntegral}
\lean{ComputableAnalysis.ExponentialComputations.integralLog_hasIntegral}
\lean{ComputableAnalysis.ExponentialComputations.integralLog_equiv_log}
\lean{ComputableAnalysis.ExponentialComputations.integralLog_congr}
\leanok
For every positive valid represented endpoint $x$, the rectangle-sum
computation has the supplied reciprocal integral and
\[
 L(x)\simeq\log x.
\]
The computation and the identity are representation invariant.
\end{theorem}

This uses tight rectangles as a proved sufficient construction for this
integrand. It does not require every future integral computation to use that
criterion. A numerical candidate, or a record that merely asserts validity,
is not being promoted to an integral by definition.

\section{Exponential as the inverse of the integral}\label{sec:inverse-integral-exponential}

Combining the independently constructed integral agreement with the global
inverse laws proves
\[
 L(\exp a)\simeq a,\qquad \exp(L(x))\simeq x.
\]
Thus $L$ is onto the represented real axis and has a unique positive inverse
value at every valid represented real input. This is an exact characterization
of exponential by integration, with its existence and uniqueness proved.

\lean{ComputableAnalysis.ExponentialComputations.integralLog_exp}
\lean{ComputableAnalysis.ExponentialComputations.exp_integralLog}
\lean{ComputableAnalysis.ExponentialComputations.exp_unique_of_integralLog}
\leanok

A future independent algorithm that inverts $L$ can join the API by proving
its positive-domain evidence and the equation $L(f(a))\simeq a$. The checked
uniqueness theorem then proves $f(a)\simeq\exp a$. There is currently no
separate exponential evaluator that runs bisection using only integral
samples. The two-sided inverse characterization above is complete; that
additional numerical backend is a separate construction.

\section{An API for interchangeable computations}\label{sec:exponential-computation-api}

Proving all pairwise comparisons would duplicate work. Instead, each
exponential implementation registers one already proved agreement with
$E$. The library derives agreement between any two implementations and
invariance under equivalent input representations. The current registrations
are the power series, the literal compound-interest limit, and the scalar
linear ODE. The bounded represented Taylor construction has its own checked
bridge to the same reference. On real inputs each implementation agrees with
$\exp$.

\lean{ComputableAnalysis.ExponentialComputations.Implementation.equivalent}
\lean{ComputableAnalysis.ExponentialComputations.Implementation.congr}
\lean{ComputableAnalysis.ExponentialComputations.Implementation.realEval_agrees}
\lean{ComputableAnalysis.ExponentialComputations.linearODE}
\leanok

Logarithm computations use the same pattern. Continuation and reciprocal
rectangles register proved agreement with the global positive logarithm.
There is no need to reproduce every logarithm proof whenever a new exponential
method is introduced, or conversely.

\lean{ComputableAnalysis.ExponentialComputations.LogImplementation.equivalent}
\lean{ComputableAnalysis.ExponentialComputations.LogImplementation.congr}
\lean{ComputableAnalysis.ExponentialComputations.LogImplementation.exp_eval}
\lean{ComputableAnalysis.ExponentialComputations.real_computation_agrees_of_integral_inverse}
\leanok

Finally, $e=\exp(1)$ agrees with the repository's earlier power-series and
compound-interest constants. Those compatibility theorems are separate from
the new core's axiom audit: they inherit older native arithmetic proof
certificates. The core chapter endpoints are audited against only the standard
Lean axioms. The older named rational computation \texttt{logTwoSeries}
still needs its own bridge to this global API; the reciprocal computation
constructed here already agrees with $\log 2$ by the global theorem.

\lean{ComputableAnalysis.ExponentialComputations.exp_one_equiv_ePowerSeries}
\lean{ComputableAnalysis.ExponentialComputations.exp_one_equiv_eCompoundInterest}
\leanok

\section{Complex exponential and logarithm branches}\label{sec:complex-exp-log}

The exponential API starts on the complex plane; its real version is a
restriction. Every input is a valid represented complex number, computed
by shrinking rational rectangles. The same power series, compound-interest
limit, and scalar ODE computations therefore work before choosing a real axis.

\begin{theorem}[Entire complex exponential]
\lean{ComputableAnalysis.ExponentialComputations.Complex.exp}
\lean{ComputableAnalysis.ExponentialComputations.Complex.expHolomorphic}
\lean{ComputableAnalysis.ExponentialComputations.Complex.exp_derivative}
\lean{ComputableAnalysis.ExponentialComputations.Complex.exp_add}
\lean{ComputableAnalysis.ExponentialComputations.Complex.exp_nonzero}
\lean{ComputableAnalysis.ExponentialComputations.Complex.exp_real}
\leanok
The represented complex exponential is holomorphic everywhere, has derivative
$E'(z)=E(z)$, never vanishes, and satisfies $E(z+w)\simeq E(z)E(w)$.
On the real axis it agrees as a full complex value with the embedding of
the real exponential.
\end{theorem}

For logarithm, a normalization fixes a branch. Choose a nonzero center $c$
and a represented value $a$ with $E(a)\simeq c$. Near that center compute
\[
 L_{c,a}(z)=a+\sum_{n=1}^{\infty}\frac{(-1)^{n+1}}{n}
                    \left(\frac{z}{c}-1\right)^n.
\]
The actual open chart consists of inputs for which
$\|z/c-1\|_{\square}<1/32$, with a constructive strict-bound witness.
The small radius is a quantitative choice for the current estimates.
The seed supplies only a center normalization; holomorphicity, the derivative,
and the inverse identity are proved from this computation.

\begin{theorem}[Normalized complex logarithm]
\lean{ComputableAnalysis.ExponentialComputations.Complex.LogSeed}
\lean{ComputableAnalysis.ExponentialComputations.Complex.LogSeed.function}
\lean{ComputableAnalysis.ExponentialComputations.Complex.LogSeed.holomorphic}
\lean{ComputableAnalysis.ExponentialComputations.Complex.LogSeed.initial}
\lean{ComputableAnalysis.ExponentialComputations.Complex.LogSeed.derivative}
\lean{ComputableAnalysis.ExponentialComputations.Complex.LogSeed.exponential_eval}
\lean{ComputableAnalysis.ExponentialComputations.Complex.LogSeed.congr}
\lean{ComputableAnalysis.ExponentialComputations.Complex.LogSeed.log_exp}
\leanok
On the actual chart, $L_{c,a}(c)\simeq a$,
$L'_{c,a}(z)\simeq 1/z$, and $E(L_{c,a}(z))\simeq z$.
Transporting equivalent centers, seeds, and inputs preserves its values.
Conversely, $L_{c,a}(E(w))\simeq w$ when $E(w)$ is in this chart and
$\|w-a\|_{\square}\leq 1$; proximity to the chosen seed fixes the branch.
\end{theorem}

To continue a branch, move to a point of its chart and use the computed value
there as the next seed. The finite chain records genuine chart membership;
convexity certifies every point of each affine edge. An explicit open radius
at each transition gives agreement of the old and new germs. Admissible
subdivision preserves the terminal value. The API carries the normalization
through the chain instead of choosing a fresh branch at every step.

\lean{ComputableAnalysis.ExponentialComputations.Complex.LogSeed.recenter}
\lean{ComputableAnalysis.ExponentialComputations.Complex.LogSeed.recenter_agrees}
\lean{ComputableAnalysis.ExponentialComputations.Complex.LogSeed.continueAlong}
\lean{ComputableAnalysis.RiemannHilbert.LogarithmContinuation.edge_segment}
\lean{ComputableAnalysis.RiemannHilbert.LogarithmContinuation.edge_germ}
\lean{ComputableAnalysis.RiemannHilbert.LogarithmContinuation.subdivision}
\leanok

\begin{theorem}[Extension of the positive real logarithm]
\lean{ComputableAnalysis.ExponentialComputations.Complex.realLogarithmAt}
\lean{ComputableAnalysis.ExponentialComputations.Complex.realLogarithmAt_initial}
\lean{ComputableAnalysis.ExponentialComputations.Complex.realLogarithmAt_agrees}
\leanok
At every positive represented real center $x$, choose the real seed
$a=\log(x)$. Throughout the positive real part of that chart the complex
branch agrees with the full embedding of $\log(y)$. Thus the real reciprocal
integral and its inverse laws from the previous sections retain their meaning
inside the complex branch API.
\end{theorem}

There is no single holomorphic logarithm on the entire punctured plane.
The executable route selection constructs one logarithm and a local germ at
any nonzero input, but independently selected values are not asserted to form
a globally continuous function. A principal branch would need its own cut
domain and agreement proof. Likewise, a complex path-integral implementation
of logarithm must supply actual complex integral evidence; the derivative
identity alone is not that evidence.

\lean{ComputableAnalysis.ExponentialComputations.Complex.logarithmAt}
\leanok

\begin{theorem}[Branch agreement and periods]
\lean{ComputableAnalysis.ExponentialComputations.Complex.LogSeed.overlap_difference}
\lean{ComputableAnalysis.ExponentialComputations.Complex.LogSeed.overlap_agrees}
\lean{ComputableAnalysis.ExponentialComputations.Complex.period}
\lean{ComputableAnalysis.ExponentialComputations.Complex.exp_period}
\lean{ComputableAnalysis.ExponentialComputations.Complex.exp_translate_period}
\lean{ComputableAnalysis.ExponentialComputations.Complex.exp_kernel}
\lean{ComputableAnalysis.ExponentialComputations.Complex.exp_fiber}
\lean{ComputableAnalysis.ExponentialComputations.Complex.LogSeed.overlap_period}
\leanok
Equal exponential values differ by $2\pi i k$ for some integer $k$, and every
such translation preserves the exponential. On a nonempty intersection of
two of these convex logarithm charts, one fixed integer $k$ describes the
difference everywhere. Agreement at one point forces agreement throughout
the overlap, by a proved finite telescoping argument.
\end{theorem}

Here $2\pi i$ is four times the imaginary geometric half-pi in the computable
foundation. The period classification reuses the existing checked nome-kernel
classification after rescaling by this nonzero period. Its current import
closure contains auxiliary modular-function theory; factoring that dependency
into a smaller general exponential-period module is a later organizational task.
